% ============================================================================
\section{Discussion}
\label{sec:discussion}
% ============================================================================
\subsection{Effect of the Two Layers and the Two Time Axes}
The system paid more employees correctly on the payroll date than the code-based update in four of the five companies and matched or exceeded that update in all five with the completeness gate. The advantage comes from timing, since an approved update enters the payroll on the issue date of the document while a change in program code waits for a developer. The comparison favors the code-based update, which is always correct once it arrives and ends every company at 1.000. The system carries its silent errors to the end of the window. On the other side, the comparison credits the system with applying an approved update on the issue date, and the size of its advantage depends on the 26-day delay that represents the fastest observed code change.

The bitemporal records make this timing safe for retroactive rules. A correct retroactive update that arrives after a payroll has run produces the exact correction for the months already paid. The payroll as computed before the rule stays reproducible for the employer's records and for a tax audit. Neither property required program code in any of the 15 late-upload scenarios.

The two layers let every company start from the same validated regulation layer while keeping its own policies. A sector-specific regulation such as the DTP facility of PMK 10/2025 changed the payslips of eligible companies without a company-specific rule because the engine checks its conditions against the profile of each company.

\subsection{Reliability of LLM-Drafted Updates}
The drafted updates were most reliable on government regulations. All parameter changes and all changes of calculation method were correct, as were 35 of the 39 proposals on prospective regulations, and no distractor changed a payslip. Company policies were harder: the system abstained on all six real labor agreements and made silent errors on 11 of the 30 prospective policy units. Five of these errors arose only because the amended policy was named differently from the reference.

The gates mattered most against the selection rule of the closest work~\cite{goganikhiabani2025maintaining}, a majority vote without gates. This majority vote made 24 silent errors on the prospective set, whereas the gated selections made 13 to 15. The behavioral consensus added no accuracy over taking the first draft that passes the gates and serves mainly to decide when the system abstains. The gate on payroll inputs (G6) rejected the most drafts. A rule that reads an attribute missing from the payroll data passes the tests on generated cases because these cases supply every attribute. G6 stops such rules before they reach the engineer.

\subsection{Role of the HC Software Engineer}
The system is designed for an HC software engineer who approves proposals and is not meant to run unattended. Without a reviewer, accuracy fell to between 0.000 and 0.566 because the abstained labor agreements left the company layers empty. The system decided 113 of the 138 documents when an engineer approved its proposals. For 11 of the other documents the engineer wrote the updates as structured changes with quotes instead of new program code. The reviewer condition also had the engineer approve 10 silently wrong proposals and is the worst case because it assumes that the engineer never notices an error. Some of these errors carry signals that a reader of the proposal can see, such as the note in the proposal for Company D that the update was not computationally complete. The careful-reviewer condition models an engineer who checks the review payslips. It reached 1.000 at the end of the window at Companies A and C. Its assumption that every visible error is noticed makes it an upper bound for this kind of check.

Text in a document that tries to steer the model cannot reach a payslip without the engineer's approval, and any new number that it introduces must still be quoted from the document (G2). A forged document with quotes that support a wrong rule passes these checks, so the engineer must also confirm where each document comes from.

\subsection{Payslip-Level Testing of Chained Updates}
Eight of the 48 pairs of a regulation and a real company had a different outcome in the chain than in the isolated run, and the labor agreement of Company D adds a ninth. The 2026 DTP regulation was correct whenever it was drafted against the reference knowledge base and wrong in four chains, where the facility rules it extends had been written by the system. Prior et al.\ observed a related effect for consolidated statute text: exact match fell from 50.3\% for one amendment to 20.51\% along amendment chains~\cite{prior2025risks}. A test of each update on its own therefore cannot show how a company's payslips will develop. The chained, payslip-level replay used here can serve as an acceptance test for other systems that maintain executable rules with machine-drafted updates. The gates can be strengthened in the same direction by generating review cases from the company's current knowledge base and from the months in which an update takes effect.

\subsection{Limitations}
The evaluation covers five real companies in five sectors, one country, and one tax. The payslip results therefore describe these companies and the PPh 21 rules of 2023--2026. Because each chain was run once, the result of Company D rests on one stochastic proposal for a document on which both isolated runs abstained. The prospective units are shorter than real regulations and were written for the study by a model of the same family as the second drafting model. Their difficulty may therefore differ from that of future real changes. The reference updates of these units were checked mechanically but not by the tax practitioner. The results on distractors apply to trimmed texts that keep only the opening, one article, and the closing of each distractor. The effect of the completeness gate on new documents remains to be confirmed because the gate was defined after the measured runs and evaluated by replaying their recorded drafts without a repair round. The reviewer conditions model the HC software engineer instead of observing one. The main results come from one drafting model whose training data may include the regulations of the test period.
